% ---
% 2 Tecnologias Envolvidas
% Este arquivo descreve as tecnologias utilizadas no projeto.
% Conforme as Diretrizes do PPC, deve-se justificar as escolhas tecnológicas.
% ---
\chapter{TECNOLOGIAS ENVOLVIDAS}
\label{cap:tecnologias}

Neste capítulo, devem ser descritas todas as tecnologias, ferramentas, linguagens de programação, \textit{frameworks} e bibliotecas utilizados no desenvolvimento do projeto. Conforme as Diretrizes do PPC, o objetivo é justificar as escolhas tecnológicas com base em critérios técnicos.

Não se limite a listar as ferramentas; explique o papel de cada uma no contexto da arquitetura do software e os motivos que levaram à sua adoção em detrimento de alternativas disponíveis no mercado.

\section{Linguagens de Programação e Frameworks}

Descreva as linguagens utilizadas no \textit{frontend} e \textit{backend}, bem como os \textit{frameworks} adotados. Por exemplo, justifique o uso de Python com Django para o \textit{backend} devido à sua produtividade \cite{django_docs}, ou React para o \textit{frontend} devido à sua capacidade de criar interfaces reativas.

\section{Banco de Dados}

O Sistema Gerenciador de Banco de Dados (SGBD) escolhido para o desenvolvimento do sistema de Licença Capacitação foi o \textbf{MySQL}, em sua versão estável mais recente. A escolha se justifica principalmente pela necessidade de \textbf{integridade relacional consistente}, uma vez que o domínio do problema é regido por regras legais rígidas oriundas da Lei nº 8.112/1990, do Decreto nº 9.991/2019 e da Instrução Normativa nº 21/2021, que precisam ser refletidas em nível de banco de dados por meio de restrições de integridade (chaves estrangeiras, checagens de unicidade e validações de domínio).

Entre os fatores que motivaram a escolha, destacam-se:

\begin{itemize}
\item \textbf{Suporte a constraints e integridade referencial} (chaves estrangeiras, \texttt{UNIQUE}, \texttt{CHECK}), essencial para impedir, por exemplo, o cadastro de parcelas de afastamento com sobreposição de datas;
\item \textbf{Ampla adoção no mercado e em sistemas do setor público}, o que facilita manutenção futura, contratação de suporte e integração com outras soluções institucionais;
\item \textbf{Boa performance em operações de leitura e escrita}, adequada ao volume de dados esperado para o sistema (cadastro de servidores, solicitações e parcelas de afastamento);
\item \textbf{Gratuidade (Community Edition) e ampla documentação}, adequado ao contexto de um projeto acadêmico sem custos de licenciamento;
\item \textbf{Curva de aprendizado reduzida e ferramental maduro} (MySQL Workbench, phpMyAdmin), o que facilita a modelagem, administração e documentação técnica exigida por este relatório.
\end{itemize}

Optou-se por um modelo estritamente relacional, em detrimento de soluções NoSQL, dado que o domínio da aplicação exige consistência transacional (utilizando o mecanismo de armazenamento InnoDB, que garante conformidade com as propriedades ACID) e relacionamentos bem definidos entre entidades como servidor, solicitação, tramitação e parcela — características para as quais os SGBDs relacionais são mais adequados. Para mais detalhes sobre o MySQL, consulte a documentação oficial \cite{mysql_manual}.

\section{Ferramentas de Desenvolvimento e Controle de Versão}

Para o desenvolvimento do sistema, foram utilizadas as seguintes ferramentas de apoio:

\begin{itemize}
    \item \textbf{IDE/Editor:} \textbf{Visual Studio Code (VS Code)}, utilizado por toda a equipe para o desenvolvimento do backend, frontend e scripts de banco de dados, devido à sua leveza, extensibilidade e suporte nativo a diversas linguagens;
    \item \textbf{Controle de Versão:} \textbf{Git}, associado ao \textbf{GitHub}\footnote{GitHub: \url{https://github.com}} como plataforma de hospedagem remota, permitindo o versionamento colaborativo do código-fonte e o acompanhamento do histórico de alterações entre os membros da equipe;
    \item \textbf{Ferramentas de Modelagem:} \textbf{draw.io}, utilizada na elaboração do Diagrama Entidade-Relacionamento (DER) em notação clássica (retângulos para entidades, elipses para atributos, elipses sublinhadas para chaves primárias e losangos para relacionamentos), apresentado na Seção \ref{cap:modelagem};
    \item \textbf{Validação de Requisitos:} \textbf{Google Forms}, utilizado para a coleta de \textit{feedback} junto às partes interessadas (\textit{stakeholders}) durante a fase de validação das regras de negócio do módulo de banco de dados.
\end{itemize}
